%% LyX 2.5.1 created this file.  For more info, see https://www.lyx.org/.
%% Do not edit unless you really know what you are doing.
\documentclass[12pt,english]{article}
\usepackage[utf8]{inputenc}
\synctex=-1
\usepackage{color}
\usepackage{babel}
\usepackage{amsmath}
\usepackage{amssymb}
\usepackage[bookmarks=false,
 breaklinks=false,pdfborder={0 0 1},backref=false,colorlinks=true]
 {hyperref}
\hypersetup{
 allcolors=blue}

\makeatletter
\@ifundefined{date}{}{\date{}}
%%%%%%%%%%%%%%%%%%%%%%%%%%%%%% User specified LaTeX commands.
\PassOptionsToPackage{english}{babel}
\usepackage{babel}
\title{Draft Appendix: Relativistic Hamiltonian and Scalar Wave Equations}
\author{}



\makeatother

\begin{document}
\maketitle% Included from main.tex after \appendix.
% This file intentionally does not redefine the document class or packages.

\subsection{A Note on the Energy Eigenvalue in the Corrected Schrödinger Equation}

\label{app:corrected_schrodinger_eigenvalue}

The relativistically corrected Schrödinger equation is commonly written
as

In particular, the physical time dependence of a stationary relativistic
state remains 
\begin{equation}
\Psi(\mathbf{r},t)=\psi(\mathbf{r})\exp\left(-\frac{iE_{\mathrm{tot}}t}{\hbar}\right).
\end{equation}
Thus, $E_{\mathrm{tot}}$ governs the temporal phase, whereas $p^{2}_{0}/(2\gamma m_{0})$
is the eigenvalue of the equivalent stationary Schrödinger-like equation.

\section{Relativistic Hamiltonian and Scalar Wave Equations}

\label{app:relativistic_hamiltonian}

This appendix gives the derivation underlying the two scalar refractive-index
models used in Section~\ref{sec:schrodinger_helmholtz}. A complete
relativistic description begins with the Dirac equation. When spin
and magnetic interactions are not explicitly required, elastic electron
diffraction can instead be described by the positive-energy, spin-independent
relativistic Hamiltonian \cite{fujiwara1961relativistic,kirkland2010advanced}.

\subsection{Minimal Coupling and the Stationary Equation}

The time-dependent wave equation is 
\begin{equation}
i\hbar\frac{\partial\Psi(\mathbf{r},t)}{\partial t}=\hat{H}\Psi(\mathbf{r},t),\label{eq:app_time_dependent_wave_equation}
\end{equation}
with the minimal-coupling Hamiltonian 
\begin{equation}
\hat{H}=\sqrt{c^{2}\left(\hat{\mathbf{p}}-q_{e}\mathbf{A}/c\right)^{2}+m^{2}_{0}c^{4}}+q_{e}V(\mathbf{r}),\label{eq:app_minimal_coupling_hamiltonian}
\end{equation}
where $\hat{\mathbf{p}}=-i\hbar\nabla$, $\mathbf{A}$ is the magnetic
vector potential, and $V(\mathbf{r})$ is the specimen electrostatic
potential in volts. We use $e>0$ and $q_{e}=-e$, so the electron
potential energy is 
\begin{equation}
E_{\mathrm{pot}}(\mathbf{r})=q_{e}V(\mathbf{r})=-eV(\mathbf{r}).\label{eq:app_electron_potential_energy}
\end{equation}

This notation distinguishes the electrostatic potential $V$ from
the energy $E_{\mathrm{pot}}$.

Neglecting magnetic interactions and introducing the rest energy $E_{0}=m_{0}c^{2}$
gives 
\begin{equation}
\hat{H}=\sqrt{c^{2}\hat{\mathbf{p}}^{\,2}+m^{2}_{0}c^{4}}+E_{\mathrm{pot}}(\mathbf{r})=\sqrt{c^{2}\hat{\mathbf{p}}^{\,2}+E^{2}_{0}}+E_{\mathrm{pot}}(\mathbf{r})\label{eq:app_relativistic_hamiltonian}
\end{equation}

We then look for stationnary states $\Psi(\mathbf{r},t)$ that separate
the temporal and spatial contributions : 
\begin{equation}
\Psi(\mathbf{r},t)=\psi(\mathbf{r})\exp\left(-\frac{iE_{\mathrm{tot}}t}{\hbar}\right),\label{eq:app_stationary_ansatz}
\end{equation}

The stationary equation is therefore 
\begin{equation}
\hat{H}\psi=E_{\mathrm{tot}}\psi.\label{eq:app_stationary_equation}
\end{equation}

$E_{\mathrm{tot}}$ is the conserved total relativistic energy, including
the rest energy. As we look for elastic interaction this energy is
conserved. In particular, in vacuum, 
\begin{equation}
E_{\mathrm{tot}}=\gamma m_{0}c^{2}=\sqrt{c^{2}p^{2}_{0}+E^{2}_{0}},\qquad p_{0}=hk_{0}=\frac{h}{\lambda},\label{eq:app_vacuum_energy_momentum}
\end{equation}
. If the beam has accelerated through a voltage magnitude $U_{\mathrm{acc}}>0$,
the total energy is the sum of the rest and kinetic energies :

and $E_{\mathrm{kin}}=eU_{\mathrm{acc}}=E_{\mathrm{tot}}-E_{0}=(\gamma-1)E_{0}$.

Equation~\eqref{eq:app_vacuum_energy_momentum} can be used to compute
the wave length : $c^{2}p^{2}_{0}=h^{2}c^{2}/\lambda^{2}=E^{2}_{\mathrm{tot}}-E^{2}_{0}=(E_{\mathrm{kin}}+E_{0})^{2}-E^{2}_{0}$

\subsection{Klein--Gordon and Corrected Schrödinger Forms}

Rearranging Eq.~\eqref{eq:app_stationary_equation} and squaring
gives 
\begin{equation}
\left[c^{2}\hat{\mathbf{p}}^{\,2}+E^{2}_{0}\right]\psi=\left[E_{\mathrm{tot}}-E_{\mathrm{pot}}(\mathbf{r})\right]^{2}\psi.\label{eq:app_kg_before_linearisation}
\end{equation}
A local relattivistic momentum and a local wavelength can be defined
: 
\begin{equation}
c^{2}\hat{\mathbf{p}}^{\,2}=\left[E_{\mathrm{tot}}-E_{\mathrm{pot}}(\mathbf{r})\right]^{2}-E^{2}_{0}=h^{2}c^{2}\mathbf{k}^{2}\label{eq:app_kg_before_linearisation2}
\end{equation}

Using $\hat{\mathbf{p}}^{\,2}=-\hbar^{2}\nabla^{2}$ and $p_{0}=hk_{0}$,
equation ~\eqref{eq:app_vacuum_energy_momentum} becomes 
\begin{equation}
\nabla^{2}\psi(\mathbf{r})+(2\pi k_{0})^{2}n^{2}_{\mathrm{KG}}(\mathbf{r})\psi(\mathbf{r})=0,\label{eq:app_kg_helmholtz_form}
\end{equation}
where 
\begin{equation}
n^{2}_{\mathrm{KG}}(\mathbf{r})=\frac{\left[E_{\mathrm{tot}}-E_{\mathrm{pot}}(\mathbf{r})\right]^{2}-E^{2}_{0}}{E^{2}_{\mathrm{tot}}-E^{2}_{0}}=\frac{p(\mathbf{r})}{p_{0}}=\frac{k(\mathbf{r})}{k_{0}}.\label{eq:app_kg_refractive_index}
\end{equation}
So, $n_{\mathrm{KG}}$ is the ratio of the local relativistic momentum
$p(r)$ to the vacuum momentum $p_{0}$ . This connects the Klein--Gordon
form to the electron-optical refractive-index concepts of Glaser and
Ehrenberg--Siday \cite{Glaser1933,Ehrenberg-Siday1949}.

For $|E_{\mathrm{pot}}|\ll E_{\mathrm{tot}}$, as it is the case for
electron-beam interaction, 
\begin{equation}
\left[E_{\mathrm{tot}}-E_{\mathrm{pot}}(\mathbf{r})\right]^{2}\simeq E^{2}_{\mathrm{tot}}-2E_{\mathrm{tot}}E_{\mathrm{pot}}(\mathbf{r}).\label{eq:app_linearised_energy}
\end{equation}
Since $E_{\mathrm{tot}}/c^{2}=\gamma m_{0}$, Eq.~\eqref{eq:app_kg_before_linearisation}
then reduces to 
\begin{equation}
\nabla^{2}\psi(\mathbf{r})+\left[(2\pi k_{0})^{2}-\frac{2\gamma m_{0}}{\hbar^{2}}E_{\mathrm{pot}}(\mathbf{r})\right]\psi(\mathbf{r})=0.\label{eq:app_corrected_schrodinger}
\end{equation}
or equivalently ,

\begin{equation}
\left[-\frac{\hbar^{2}}{2\gamma m_{0}}\nabla^{2}+E_{\mathrm{pot}}(\mathbf{r})\right]\psi(\mathbf{r})=\frac{p^{2}_{0}}{2\gamma m_{0}}\psi(\mathbf{r}).\label{eq:app_schrodinger_energy_form}
\end{equation}

Equation~\eqref{eq:app_schrodinger_energy_form} is therefore a stationary
Schrödinger-like reformulation of the relativistic wave equation,
rather than the original relativistic energy-eigenvalue equation.
This is the relativistically corrected Schrödinger form as deduced
by Fujiwara. The right-hand side resembles the nonrelativistic kinetic
energy $p^{2}_{0}/(2m)$, with the rest mass replaced by the relativistic
mass $\gamma m_{0}$. It should not, however, be identified with either
the total relativistic energy $E_{\mathrm{tot}}$ or, in general,
the kinetic energy $E_{\mathrm{kin}}=E_{\mathrm{tot}}-E_{0}$.

Indeed, using 
\begin{equation}
\gamma m_{0}=\frac{E_{\mathrm{tot}}}{c^{2}},\qquad p^{2}_{0}c^{2}=E^{2}_{\mathrm{tot}}-E^{2}_{0},
\end{equation}
the eigenvalue in Eq.~\eqref{eq:app_schrodinger_energy_form} becomes
\begin{equation}
\frac{p^{2}_{0}}{2\gamma m_{0}}=\frac{E^{2}_{\mathrm{tot}}-E^{2}_{0}}{2E_{\mathrm{tot}}}=E_{\mathrm{kin}}\frac{E_{\mathrm{tot}}+E_{0}}{2E_{\mathrm{tot}}}.\label{eq:app_effective_schrodinger_eigenvalue}
\end{equation}

It approaches $E_{\mathrm{kin}}$ only in the nonrelativistic limit
$E_{\mathrm{kin}}\ll E_{0}$, for which $E_{\mathrm{tot}}\simeq E_{0}$.

Equation xxxx-mettrelabonneREF indicates that the energy appearing
in the temporal dependence of the wave function is not $E_{0\mathrm{kin}}$
but $E_{\mathrm{tot}}$. The factor $1/2$ arises when the scalar
relativistic wave equation is linearised in the specimen potential
and divided by $2E_{\mathrm{tot}}$.

\subsection{Interaction Constant and Paraxial Limit}

The conventional TEM interaction constant is 
\begin{equation}
\sigma=\frac{2\pi\gamma m_{0}e}{h^{2}k_{0}}=\frac{2\pi\gamma m_{0}e\lambda}{h^{2}}=2\pi ek_{0}\frac{\gamma m_{0}c^{2}}{h^{2}k^{2}_{0}c^{2}}=\frac{2\pi eE_{\mathrm{tot}}}{hc\sqrt{E^{2}_{\mathrm{tot}}-E^{2}_{0}}}.\label{eq:app_sigma_definition}
\end{equation}

The potential term in Eq.~\eqref{eq:app_corrected_schrodinger} can
therefore be written as 
\begin{equation}
\frac{2\gamma m_{0}e}{\hbar^{2}}V(\mathbf{r})=4\pi k_{0}\sigma V(\mathbf{r}),\label{eq:app_potential_term_sigma}
\end{equation}
giving 
\begin{equation}
\nabla^{2}\psi(\mathbf{r})+\left[(2\pi k_{0})^{2}+4\pi k_{0}\sigma V(\mathbf{r})\right]\psi(\mathbf{r})=0.\label{eq:app_corrected_schrodinger_sigma}
\end{equation}
Factoring out the carrier wave, 
\begin{equation}
\psi(\mathbf{r})=\exp(2\pi ik_{0}z)\tilde{\psi}(\mathbf{r}),\label{eq:app_carrier_wave}
\end{equation}
and neglecting $\partial^{2}\tilde{\psi}/\partial z^{2}$ gives 
\begin{equation}
\frac{\partial\tilde{\psi}}{\partial z}=\frac{i}{4\pi k_{0}}\nabla^{2}_{\perp}\tilde{\psi}+i\sigma V(\mathbf{r})\tilde{\psi}.\label{eq:app_paraxial_sigma}
\end{equation}
For a slice with projected electrostatic potential 
\begin{equation}
V_{j}(x,y)=\int^{z_{j}+\Delta z}_{z_{j}}V(x,y,z)\,dz,\label{eq:app_projected_potential}
\end{equation}
the phase-grating transmission function is 
\begin{equation}
t_{j}(x,y)=\exp\left[i\sigma V_{j}(x,y)\right],\label{eq:app_transmission_function}
\end{equation}
which is the expression used in the main text.

 \bibliographystyle{elsarticle-num}
\bibliography{sample}
 
\end{document}
